\documentclass[11pt,a4paper]{article}
\usepackage{lmodern}
\usepackage[utf8]{inputenc}
\usepackage[T1]{fontenc}
\usepackage{amsmath}
\usepackage{amsfonts}
\usepackage{amssymb}
\usepackage{graphicx}
\usepackage[french]{babel}
\usepackage[notitlepage,fancysections,pagenumber]{sauveterre}

\usepackage{enumitem}

\usepackage{indentfirst}
\usepackage{lipsum}

\usepackage{cases}
\usepackage{multicol}

%formatting titles
\usepackage{titlesec}
\titleformat*{\section}{\centering\Large\bfseries}
\titleformat*{\subsection}{\centering\large\bfseries}

%formatting date
\usepackage[ddmmyyyy]{datetime}
\renewcommand{\dateseparator}{/}

%\usepackage{geometry}
%\geometry{
% a4paper,
% total={170mm,257mm},
% left=20mm,
% top=20mm,
% }

\usepackage{pgf,tikz}
\usepackage{mathrsfs}
\usetikzlibrary{arrows}

\usepackage[export]{adjustbox}

\renewcommand{\notitlepageupperline}{
	\noindent%
	\begin{minipage}{\textwidth}
		\centering
		\begin{minipage}{0.33\textwidth}
			\raisebox{-0.5\height}
			{\color{bleu-sauveterre}%
				\includegraphics[height=0.8\hauteurlogovert]
				{\polylogovert}
				}
		\end{minipage}%
		\begin{minipage}{0.33\textwidth}
			\color{bleu-sauveterre}%
			\centering\polydatesave
			\polyauthorsave
			\hspace{3em}
		\end{minipage}%
		\begin{minipage}{0.33\textwidth}
			\color{bleu-sauveterre}%
			\flushright\raisebox{-0.5\height}
			{\includegraphics[height=0.8\hauteurlogovert]{logo-federal}
			}
		\end{minipage}
	\end{minipage}
}

%\usepackage{fancyhdr}
%\usepackage{lastpage}
%\pagestyle{fancy}
%\fancyhf{}
%%headers and footers
%\lhead{Centre Équestre de Sauveterre}
%\chead{Règlement intérieur}
%\rhead{Version du \today}
%\rfoot{Page \thepage \ sur \pageref{LastPage}}

%\fancypagestyle{firstpage}{%
%  \fancyhf{}
%  \renewcommand\headrulewidth{0pt}
%  \fancyfoot[R]{Page \thepage \ sur \pageref{LastPage}}
%}
%
%\title{\vspace{-50pt}
%\includegraphics[width=0.4\textwidth , left]{logo2-v2.eps}~
%\\[8pt]
%{\huge \textsc{Règlement intérieur}}
%}

\title{Conditions Générales applicables aux Balades}
\author{}
\date{}

%enumerate environment formatting
\renewcommand{\labelenumi}{\textbf{\arabic{enumi}.}}
\renewcommand{\labelenumii}{\textbf{\arabic{enumi}.\arabic{enumii}.}}
\renewcommand{\labelenumiii}{\textbf{\arabic{enumi}.\arabic{enumii}.\arabic{enumiii}.}}

%change sec and subsec numbering
\renewcommand{\thesection}{Section \Roman{section}.}
\renewcommand{\thesubsection}{Article \arabic{subsection} :}

\renewcommand{\labelenumi}{\arabic{enumi}.}


\renewcommand*\contentsname{Plan}


\usepackage{framed}

%\setlength{\parindent}{0pt}%remove indentation

\begin{document}

\maketitle

%\thispagestyle{firstpage}
%\vspace{-30pt}

\subsection*{Balades à poneys}
Les tours de poneys en main se font sous l'entière responsabilité des parents, de la famille, ou de leur délégué. \par
L'accompagnant familial doit disposer de toutes ses facultés physiques pour assurer la sécurité du cavalier ou de la cavalière lors de la prise en main. Tout handicap reconnu (articulations foulées ou cassées, etc...) ou situation ponctuelle (femme enceinte, etc...) interdit la prise en main pour des raisons évidentes de sécurité.

Les balades en main se font uniquement au pas, à raison d'un enfant seulement par poney, ce dernier devant obligatoirement porter une bombe. Le non respect de ces conditions peut entrainer une exclusion du Centre Équestre comme prévu dans l'article 8.

\subsection*{Balades à cheval}
Les balades à cheval se font uniquement sur réservation. Les réservations pour les balades se font en personne à l'accueil du Centre Équestre, par téléphone, ou par courrier électronique.

Toute balade à cheval est encadrée par le nombre d'enseignant défini par la direction, en accord avec les règlements et autorisations afférents aux diplômes et assurances.

On distingue, en balade, deux groupes distincts : les \og débutants \fg{}, \textit{i.e.} toute personne non expérimentée montant à cheval pour le première fois ou seulement occasionnellement, sans niveau fédéral, et les \og confirmés \fg{} qui sont les cavaliers possédant un niveau équestre (du point de vue fédéral) et étant en capacités de trotter et galoper en toute aisance en extérieur. Généralement, il est exigé des cavaliers qu'ils possèdent au minimum un Galop 3 pour galoper en extérieur. Tout cavalier ne répondant pas à ces critères est \textit{de facto} considéré comme débutant. Le groupe constitué des cavaliers débutants fera la balade \textit{\textbf{au pas uniquement}}. Une phase au trot est toutefois possible, si et seulement si l'enseignant responsable du groupe estime que cela ne présente pas de risques pour la sécurité du groupe. Lui seul est apte à prendre cette décision, et elle ne saurait être contestée.

Les enseignants ayant pour mission première de garantir la sécurité du groupe, ils ont la possibilité de refuser un ou plusieurs cavalier dans le groupe confirmé s'ils estiment que trotter ou galoper présente un risque pour leur sécurité. \textit{\textbf{La responsabilité des enseignants ne saurait être engagée en cas de chute d'un cavalier n'ayant pas respecté cette décision}}. De plus, aucun remboursement ne sera effectué pour un cavalier auquel aurait été refusé le galop.

\subsection*{Remboursement de la prestation}
Un faux-départ n'est pas remboursé. Aussi, toute balade commencée est due. Une balade est considérée comme \og commencée \fg{} à partir du moment où le cavalier est à cheval, et le groupe à l'extérieur du manège. De même, une balade à poney n'est plus remboursable à partir du moment où l'enfant à poney entre dans la forêt.

Dans les autres cas, la direction est la seule à décider d'un éventuel remboursement, par le biais du moyen de paiement utilisé par le règlement ou sous forme de \og carte cadeau \fg{} valable pour une autre prestation.




\end{document}
